% mazzi: 1
\chapter{La definizione di impianto sportivo e il servizio pubblico}

Il corso muove dalle definizioni e dai concetti di base --- che cosa sia ``sport'', che cosa sia
``impianto sportivo'', quali norme se ne occupino --- per arrivare al punto che regge tutta la
materia: la gestione di un impianto sportivo di proprietà pubblica è un \textbf{servizio
pubblico}, e resta tale qualunque sia la sua rilevanza economica e qualunque forma di
affidamento venga scelta.

% ---------------------------------------------------------------------------
\section{La Carta Europea dello Sport (1992)}

\textbf{Carta Europea dello Sport (1992)}: atto del legislatore comunitario che affronta il tema
degli impianti in modo ampio, declinando l'attività sportiva non necessariamente in senso
agonistico. Le sue affermazioni sono in buona parte \textbf{di principio}, ma è la sostanza che
vi sta dietro a orientare la disciplina nazionale degli impianti.

\subsection{Accesso e promozione della pratica sportiva}

\begin{itemize}
  \item \textbf{accesso non discriminatorio}: l'accesso agli impianti o alle attività sportive è
    garantito senza distinzione di sesso, razza, colore, lingua, religione, opinioni politiche o
    altra opinione, origine nazionale o sociale, appartenenza a una minoranza nazionale,
    ricchezza, nascita o qualsiasi altro status;
  \item \textbf{possibilità effettiva di praticare}: provvedimenti perché tutti i cittadini
    possano praticare lo sport, con misure supplementari --- ove necessario --- per i
    \textbf{giovani più dotati}, i \textbf{disabili} e i \textbf{gruppi svantaggiati};
  \item \textbf{accesso facilitato}: i proprietari di impianti adottano le misure necessarie per
    facilitare l'accesso dei gruppi svantaggiati, compresi coloro che soffrono di handicap fisico
    o mentale;
  \item \textbf{promozione in tutte le fasce di popolazione}: sia come divertimento, sia per
    ragioni di salute, sia per migliorare le prestazioni, offrendo impianti adeguati, programmi
    diversificati e istruttori, dirigenti o ``animatori'' qualificati;
  \item \textbf{sport sul luogo di lavoro}: la partecipazione ad attività sportive sul luogo di
    lavoro va incentivata come elemento di una politica sportiva equilibrata.
\end{itemize}

Se l'uso degli impianti pubblici deve essere garantito all'intera collettività, il ragionamento
va condotto in termini di \textbf{servizio pubblico}: la vocazione collettiva dell'impianto
pubblico è la stessa di una scuola, di un ospedale, di un teatro. La rilevanza economica della
gestione non si sottovaluta, ma la sfera e il controllo pubblicistico restano sempre in rilievo.

\subsection{La pianificazione degli impianti}

\textbf{Pianificazione globale degli impianti}: è \textbf{competenza dei poteri pubblici}, perché
la pratica sportiva dipende in parte dal numero, dalla varietà e dall'accessibilità degli
impianti.

\begin{itemize}
  \item \textbf{riferimenti della pianificazione}: esigenze nazionali, regionali e locali, e
    impianti pubblici, privati e commerciali \textbf{già esistenti};
  \item \textbf{obblighi dei responsabili}: provvedimenti per assicurare la buona gestione e la
    piena utilizzazione degli impianti, \textbf{in tutta sicurezza} (il tema della sicurezza è
    declinato dal D.M. 18 marzo 1996 e non solo);
  \item \textbf{criticità italiana}: programmazione urbanistica e impiantistica datata e scarso
    dialogo fra amministrazioni diverse $\rightarrow$ impianti dello stesso tipo vicini fra loro
    $\rightarrow$ reciproca erosione della rilevanza economica e della \textbf{sostenibilità
    economico-finanziaria};
  \item \textbf{varietà dell'offerta}: la programmazione va ragionata in termini di offerta
    sportiva complessiva, non del singolo impianto.
\end{itemize}

\subsection{La definizione di sport}

\textbf{Sport} (Carta Europea dello Sport): ``qualsiasi forma di attività fisica che, attraverso
una partecipazione organizzata o non, abbia per obiettivo l'espressione o il miglioramento della
condizione fisica e psichica, lo sviluppo delle relazioni sociali o l'ottenimento di risultati in
competizioni di tutti i livelli''.

\begin{itemize}
  \item è una definizione \textbf{più ampia} di quella corrente, che identifica lo sport con
    l'attività fisica finalizzata a una competizione --- cioè misurabile, con una classifica, e
    riconducibile a un comitato olimpico;
  \item ai fini della \textbf{programmazione pubblicistica} rileva invece anche l'attività
    organizzata \emph{o non}, e anche il solo miglioramento della condizione fisica o psichica.
\end{itemize}

% ---------------------------------------------------------------------------
\section{La definizione di sport nell'ordinamento italiano}

\textbf{Assenza di una definizione legislativa}: nella legislazione domestica non esiste una
definizione di sport. Quando la normativa italiana si riferisce allo sport, all'ordinamento
sportivo o alle politiche sportive, lo fa senza aver stabilito a priori che cosa debba
intendersi per sport.

\begin{itemize}
  \item \textbf{definizione del dizionario Garzanti}: ``l'insieme degli esercizi fisici che si
    praticano, in gruppo o individualmente, per mantenere in efficienza il corpo'';
  \item \textbf{limite della definizione}: fra le discipline sportive associate riconosciute dal
    CONI ci sono \textbf{bridge}, \textbf{dama} e \textbf{scacchi}, che non prevedono movimenti
    tesi a conservare la funzionalità del corpo $\rightarrow$ sport non equivale necessariamente a
    movimento, e il concetto di attività sportiva si estende alla \textbf{ginnastica della mente}.
\end{itemize}

\subsection{Il riconoscimento delle discipline sportive: delibera C.N. CONI 1568/2017}

\textbf{Delibera del Consiglio Nazionale CONI n.\ 1568/2017} (e successive modifiche): per la
prima volta in termini espliciti il CONI effettua una \textbf{ricognizione} delle discipline
sportive riconosciute, con elenco allegato. Ha come unico precedente una delibera del dicembre
2016.

\rubrica{Perché nasce}

\begin{itemize}
  \item le \textbf{associazioni e società sportive dilettantistiche} godono di \textbf{agevolazioni
    fiscali}, concesse a chi è riconosciuto dal CONI e iscritto al \textbf{Registro nazionale}
    delle associazioni e società sportive dilettantistiche;
  \item il CONI, quale \textbf{ente certificatore} dell'attività sportiva dilettantistica, ha
    istituito il Registro come strumento amministrativo;
  \item su sollecitazione dell'\textbf{amministrazione finanziaria} e dell'\textbf{ispettorato del
    lavoro} ha poi dovuto individuare le discipline per le quali le affiliate a enti e federazioni
    possono ottenere il riconoscimento;
  \item \textbf{effetto}: l'iscrizione al Registro, che vale come riconoscimento ai fini sportivi
    del CONI, si consegue \textbf{esclusivamente} per le discipline dell'elenco allegato.
\end{itemize}

\rubrica{Criterio utilizzato}

I tre principi ispiratori della lista originale sono le discipline sportive:
\begin{enumerate}
  \item riconosciute dal \textbf{CIO};
  \item riconosciute da \textbf{Sportaccord};
  \item presenti negli \textbf{statuti delle FSN} (federazioni sportive nazionali) \textbf{e delle
    DSA} (discipline sportive associate) approvati dalla Giunta Nazionale CONI.
\end{enumerate}

La delibera dà però atto di aver accolto parte delle richieste avanzate da FSN, DSA ed \textbf{EPS}
(enti di promozione sportiva) ``nonostante alcune integrazioni non risultino conformi'' a quei tre
principi.

\rubrica{Esito}

\textbf{385 discipline sportive riconosciute}. Il riconoscimento di alcune discipline a fronte di
altre ha generato contestazioni; fra le più note \textbf{non} comprese:

\begin{itemize}
  \item \textbf{non riconosciute}: yoga, feldenkrais, shiatsu e tutte le attività ``olistiche'';
    crossfit, TRX, pilates, zumbafitness e altre; alcune attività acquatiche e di ``fitness in
    acqua'' (attività per gestanti e neonati, idro-bike, flat, rulli in acqua); poker sportivo,
    burraco e altri giochi di carte diversi dal bridge; MMA, krav maga, paintball e simili;
  \item \textbf{riconosciute da CIO e CONI ma senza FSN o DSA}: flying disc (frisbee);
  \item \textbf{riconosciute dalla Federazione Internazionale di ginnastica ma non dal CONI}:
    sprint ostacoli e freestyle ostacoli (parkour).
\end{itemize}

% ---------------------------------------------------------------------------
\section{Le fonti legislative nazionali}

Le principali --- in realtà le uniche --- fonti normative nazionali in materia di sport:

\begin{itemize}
  \item \textbf{art.\ 117 Costituzione}: competenza concorrente delle Regioni in materia di
    ordinamento sportivo;
  \item \textbf{Legge 16 febbraio 1942 n.\ 426}: istitutiva del CONI (abrogata);
  \item \textbf{D.Lgs.\ 23 luglio 1999 n.\ 242} (\textit{decreto Melandri}), modificato dal
    \textbf{D.Lgs.\ 8 gennaio 2004 n.\ 15} (\textit{decreto Pescante}): riordino del CONI;
  \item \textbf{Legge 23 marzo 1981 n.\ 91}: norme in materia di rapporti tra società e sportivi
    professionisti;
  \item \textbf{art.\ 90 legge 289/2002} (finanziaria per il 2003): disposizioni in materia di
    attività sportiva dilettantistica.
\end{itemize}

\subsection{L'art.\ 117 Cost.\ e la competenza concorrente}

\textbf{Art.\ 117 Cost.}: l'\textbf{ordinamento sportivo} è materia di \textbf{legislazione
concorrente}; nelle materie concorrenti la potestà legislativa spetta alle Regioni, salva la
determinazione dei \textbf{principi fondamentali}, riservata alla legislazione dello Stato.

\begin{itemize}
  \item \textbf{ingresso in Costituzione}: 2001, con la riforma costituzionale del decentramento
    regionale, salutato inizialmente con favore;
  \item \textbf{meccanismo della competenza concorrente}: lo Stato definisce una normativa quadro;
    dentro la cornice così tracciata, Regioni e Province autonome di Trento e Bolzano normano in
    modo più specifico, anche in base alle connotazioni territoriali;
  \item \textbf{scelta lessicale del costituente}: non ``sport'' ma \textbf{``ordinamento
    sportivo''}, cioè un ordinamento \textbf{separato e autonomo}, non sovrano --- la sua
    sovranità cede di fronte alla supremazia dell'ordinamento statale;
  \item \textbf{difficoltà applicativa}: ``ordinamento sportivo'' è la stessa espressione con cui
    il legislatore nazionale designa l'ordinamento che fa capo al \textbf{CONI}, ente pubblico
    posto al suo vertice (decreti Melandri e Pescante).
\end{itemize}

\textbf{PROBLEMA: non esiste una legge quadro sullo sport.} Quasi tutte le Regioni hanno adottato
dopo il 2001 --- alcune con più ritocchi, qualcuna ferma a prima della riforma --- una
\textbf{legge regionale in materia di sport}, pur in assenza della legge quadro statale in materia
di ordinamento sportivo.

\subsection{Il D.P.R.\ 616/1977 sul decentramento regionale}

\textbf{Art.\ 56 D.P.R.\ 616/1977}: attribuisce alle Regioni a statuto ordinario la
\textbf{promozione di attività sportive e ricreative} e la \textbf{realizzazione dei relativi
impianti e attrezzature}.

\begin{itemize}
  \item \textbf{intesa con gli organi scolastici}: per le attività e gli impianti di interesse dei
    giovani in età scolare;
  \item \textbf{riserva al CONI}: restano ferme le attribuzioni del CONI per l'organizzazione
    delle \textbf{attività agonistiche a ogni livello} e le relative \textbf{attività
    promozionali};
  \item \textbf{consulenza tecnica}: per gli impianti e le attrezzature da essa promossi, la
    Regione si avvale della consulenza tecnica del CONI.
\end{itemize}

\textbf{Zona grigia}: se la competenza del CONI sull'organizzazione dell'attività agonistica è
chiara, è molto meno netto il confine fra le \textbf{attività promozionali connesse all'agonismo}
(CONI) e la \textbf{promozione di attività sportive e ricreative} (Regioni).

\subsection{Lo Statuto del CONI}

\textbf{Statuto CONI, art.\ 2 comma 3}: il CONI detta principi per promuovere la massima diffusione
della pratica sportiva in ogni fascia di età e di popolazione, con particolare riferimento allo
sport giovanile, sia per i normodotati sia --- di concerto con il Comitato Italiano Paralimpico
--- per i disabili, \textbf{ferme le competenze delle Regioni} e delle Province autonome di Trento
e Bolzano in materia.

\begin{itemize}
  \item il CONI è \textbf{ente pubblico} e il suo statuto è approvato con decreto della Presidenza
    del Consiglio dei Ministri;
  \item \textbf{rinvio circolare}: la norma statale fa salve le attribuzioni del CONI, lo statuto
    CONI fa salve le competenze delle Regioni --- competenze che però non sono mai state
    espressamente definite.
\end{itemize}

% ---------------------------------------------------------------------------
\section{L'impianto sportivo}

\subsection{La definizione}

\textbf{Impianto sportivo} (D.M.\ Ministero dell'Interno 18 marzo 1996): ``insieme di uno o più
spazi di attività sportiva dello stesso tipo o di tipo diverso che hanno in comune i relativi
spazi e servizi accessori, preposti allo svolgimento di manifestazioni sportive''.

\subsection{Il sistema degli impianti sportivi}

La definizione appare ovvia per lo stadio o il palazzetto, ma la nozione ampia di sport porta a
comprendere anche spazi privi di vocazione sportiva esclusiva. Il sistema comprende:

\begin{itemize}
  \item gli \textbf{stadi} e i grandi impianti per lo spettacolo sportivo;
  \item gli \textbf{impianti per le attività agonistiche} a tutti i livelli, sia pubblici sia
    privati;
  \item le \textbf{palestre} e le \textbf{piscine};
  \item gli \textbf{impianti scolastici};
  \item i \textbf{parchi urbani}, le piste ciclabili, i percorsi vita;
  \item lo \textbf{``sport senza impianti''}: fiumi, specchi d'acqua, ambiente --- definizione
    volutamente provocatoria, riferita agli sport praticabili in spazi senza vocazione sportiva
    esclusiva.
\end{itemize}

\textbf{Gestore}: responsabile della vita dell'impianto sotto il profilo tecnico, sportivo,
amministrativo, del personale, fiscale e legale.

\subsection{La natura giuridica del bene}

\textbf{Consiglio di Stato n.\ 2385/2013}: l'impianto sportivo rientra nell'ultimo capoverso
dell'\textbf{art.\ 826 c.c.}, cioè fra i beni di proprietà dei Comuni \textbf{destinati a un
pubblico servizio}.

\begin{itemize}
  \item \textbf{regime giuridico}: beni \textbf{patrimoniali indisponibili}, che ex \textbf{art.\
    828 c.c.} non possono essere sottratti alla loro destinazione;
  \item \textbf{vincolo funzionale}: coerente con la vocazione naturale del bene a essere
    impiegato in favore della collettività, per attività di interesse generale --- e la conduzione
    degli impianti sportivi rientra in tale tipologia di attività;
  \item \textbf{proprietà}: nella grande maggioranza dei casi comunale; esistono impianti di
    proprietà di province (ed ex province) o regioni.
\end{itemize}

% ---------------------------------------------------------------------------
\section{La gestione dell'impianto come servizio pubblico}

\subsection{La nozione di servizio pubblico}

\textbf{Servizio pubblico}:
\begin{itemize}
  \item \textbf{profilo dell'attività}: un complesso di operazioni funzionalmente coordinate, il
    cui prodotto è rappresentato da \textbf{utilità} poste a disposizione degli utenti per il
    soddisfacimento di \textbf{bisogni eterogenei};
  \item \textbf{profilo sostanziale}: attività \textbf{economica} che mira a soddisfare bisogni
    così largamente avvertiti da essere considerati propri di una collettività $\rightarrow$
    riconoscimento di una \textbf{funzione sociale}.
\end{itemize}

\subsection{La gestione dell'impianto sportivo}

\textbf{Gestione di un impianto sportivo}: insieme di attività volte ad assicurare il
\textbf{funzionamento} dell'impianto e l'\textbf{erogazione del servizio sportivo} che vi si
svolge.

\begin{itemize}
  \item rientra a pieno titolo nell'area dei \textbf{servizi pubblici};
  \item più in particolare nell'ambito dei \textbf{servizi alla persona} o \textbf{servizi
    sociali} --- collocazione che si ritrova poi nei singoli regolamenti comunali.
\end{itemize}

\rubrica{Conferme giurisprudenziali}

\begin{itemize}
  \item \textbf{profilo oggettivo}: ``per pubblico servizio deve intendersi un'attività economica
    esercitata per erogare prestazioni volte a soddisfare bisogni collettivi ritenuti
    indispensabili in un determinato contesto sociale, come nel caso della gestione di impianti
    sportivi comunali'' (C.d.S.\ sez.\ IV n.\ 6325/2000; C.d.S.\ sez.\ VI n.\ 1514/2001; TAR
    Lombardia Milano sez.\ III n.\ 5633/2005);
  \item \textbf{destinazione naturale del bene}: un centro sportivo strutturato in una piscina di
    proprietà comunale ``è un bene che per sua natura è destinata ad essere adibita ad un uso
    pubblico'', e l'attività inerente ha tutte le caratteristiche per essere qualificata come
    servizio pubblico (C.d.S.\ sez.\ V n.\ 4265/2008).
\end{itemize}

\subsection{La valenza pubblicistica nello Statuto CONI}

\textbf{Statuto CONI, art.\ 23 comma 1}: hanno \textbf{valenza pubblicistica} le attività delle
Federazioni sportive nazionali relative, fra l'altro, all'\textbf{utilizzazione e alla gestione
degli impianti sportivi pubblici}.

\begin{itemize}
  \item le FSN hanno natura giuridica di \textbf{associazioni riconosciute di diritto privato};
  \item per una serie di materie delegate dal CONI esercitano però \textbf{poteri
    pubblicistici}, e fra queste rientra la gestione degli impianti pubblici.
\end{itemize}

\subsection{L'apertura dell'impianto a tutti i cittadini}

\textbf{Art.\ 90 legge 289/2002, comma 24}: l'uso degli impianti sportivi in esercizio da parte
degli enti locali territoriali è \textbf{aperto a tutti i cittadini} e deve essere garantito, sulla
base di \textbf{criteri obiettivi}, a tutte le società e associazioni sportive.

\begin{itemize}
  \item è la \textbf{norma quadro} che riprende il principio di partenza della Carta Europea dello
    Sport, pur essendo contenuta in una disposizione dedicata all'attività dilettantistica;
  \item vale \textbf{indipendentemente} dalle modalità di affidamento e di gestione dell'impianto
    pubblico: il bene, quale patrimonio indisponibile gravato da vincolo funzionale, non può
    essere sottratto alla propria vocazione.
\end{itemize}
